% Part: first-order-logic
% Chapter: syntax-and-semantics
% Section: free-vars-sentences

\documentclass[../../../include/open-logic-section]{subfiles}

\begin{document}

\olfileid{fol}{syn}{fvs}

\olsection{آزاد \printtoken{P}{variable} تے \printtoken{P}{sentence}}

\begin{defn}[!!a{variable} دیاں آزاد آمدواں]
\ollabel{defn:free-occ}
!!a{formula} وچ !!a{variable} دیاں \emph{آزاد}
آمدواں اینج استقرائی طور تے تعریف کیتیاں جاندیاں نیں:
\begin{enumerate}
\item \indcase*{!A}{$!A$ ایٹمی اے}{$\indfrm$ وچ !!{variable}
  دیاں ساریاں آمدواں آزاد نیں۔}

\tagitem{prvNot}{\indcase{!A}{\lnot !B}{$\indfrm$ وچ !!{variable}
  دیاں آزاد آمدواں عین اوہو نیں جیہڑیاں $!B$ وچ آزاد نیں۔}}{}

\item \indcase{!A}{(!B \ast !C)}{$\indfrm$ وچ !!{variable}
  دیاں آزاد آمدواں $!B$ دیاں آزاد آمدواں
  تے~$!C$ دیاں آزاد آمدواں نوں ملا کے بندیاں نیں۔}

\tagitem{prvAll}{\indcase{!A}{\lforall[x][!B]}{$\indfrm$ وچ !!{variable}
  دیاں آزاد آمدواں~$!B$ دیاں ساریاں آزاد آمدواں نیں،
  سوائے~$x$ دیاں آمدواں دے۔}}{}

\tagitem{prvEx}{\indcase{!A}{\lexists[x][!B]}{$\indfrm$ وچ !!{variable}
  دیاں آزاد آمدواں~$!B$ دیاں ساریاں آزاد آمدواں نیں،
  سوائے~$x$ دیاں آمدواں دے۔}}{}
\end{enumerate}
\end{defn}

\begin{defn}[پابند متغیر]
فارمولا~$!A$ وچ !!a{variable} دی کوئی آمد
\emph{پابند} اے جے اوہ آزاد نہیں۔
\end{defn}

\begin{prob}
\olref[fol][syn][fvs]{defn:free-occ} دے طریقے تے
پابند متغیراں دیاں آمدواں دی استقرائی تعریف دیو۔
\end{prob}

\begin{defn}[دائرۂ اثر]
\iftag{prvAll}{جے $\lforall[x][!B]$، فارمولا~$!A$ وچ کسے ذیلی
  فارمولے دی اک آمد اے، تاں~$!B$ دی اوہدے مطابق آمد نوں~$!A$ وچ
  $\lforall[x]$ دی اوہدے مطابق آمد دا
  \emph{دائرۂ اثر} آکھدے آں۔ \iftag{prvEx}{$\lexists[x]$ لئی وی
  ایہو گل اے۔}{}}{جے
  $\lexists[x][!B]$، فارمولا~$!A$ وچ کسے ذیلی فارمولے
  دی اک آمد اے، تاں~$!B$ دی اوہدے مطابق آمد نوں~$!A$ وچ
  $\lexists[x]$ دی اوہدے مطابق آمد دا
  \emph{دائرۂ اثر} آکھدے آں۔}

جے $!B$ مسوّر دی کسے آمد
\iftag{prvAll}{$\lforall[x]$\iftag{prvEx}{ یا
    $\lexists[x]$}{}}{$\lexists[x]$} دا دائرۂ اثر~$!A$ وچ اے،
تاں $x$ دیاں~$!B$ وچ آزاد آمدواں
\iftag{prvAll}{$\lforall[x][!B]$\iftag{prvEx}{ تے
$\lexists[x][!B]$}{}}{$\lexists[x][!B]$} وچ پابند نیں۔
اسیں آکھدے آں کہ ایہ آمدواں مذکورہ مسوّر دی آمد
\emph{نال پابند} نیں۔
\end{defn}

\begin{ex}
ایہ فارمولا ویکھو:
\[
\lexists[\Obj v_0][\underbrace{\Atom{\Obj A^2_0}{\Obj v_0,\Obj v_1}}_{!B}] 
\]
$!B$، $\lexists[\Obj v_0]$ دا دائرۂ اثر اے۔
ایہ مسوّر $\Obj v_0$ دی آمد $!B$ وچ
پابند کردا اے، پر $\Obj v_1$ دی آمد نوں نہیں۔
ایس لئی ایس مثال وچ $\Obj v_1$ آزاد متغیر اے۔


ہُن ویکھو کہ کسے ہور پیچیدہ !!{formula}~$!A$
وچ ایہ گل کیویں کم کردی اے:
\[
\lforall[\Obj v_0][\underbrace{(\Atom{\Obj A^1_0}{\Obj v_0} \lif
    \Atom{\Obj A^2_0}{\Obj v_0, \Obj v_1})}_{!B}] \lif \lexists[\Obj
  v_1][\underbrace{(\Atom{\Obj A^2_1}{\Obj v_0, \Obj v_1} \lor \lforall[\Obj v_0][\overbrace{\lnot \Atom{\Obj A^1_1}{\Obj v_0}}^{!D}])}_{!C}]
\]
$!B$ پہلے $\lforall[\Obj v_0]$ دا دائرۂ اثر اے،
$!C$، $\lexists[\Obj v_1]$ دا دائرۂ اثر اے، تے
$!D$ دوجے $\lforall[\Obj v_0]$ دا دائرۂ اثر اے۔
پہلا $\lforall[\Obj v_0]$، $\Obj v_0$
دیاں آمدواں~$!B$ وچ پابند کردا اے؛
$\lexists[\Obj v_1]$، $\Obj v_1$ دی
آمد~$!C$ وچ پابند کردا اے؛ تے دوجا
$\lforall[\Obj v_0]$، $\Obj v_0$
دی آمد~$!D$ وچ پابند کردا اے۔
$\Obj v_1$ دی پہلی آمد تے $\Obj v_0$ دی چوتھی
آمد~$!A$ وچ آزاد نیں۔ $\Obj v_0$ دی آخری آمد
$!D$ وچ آزاد اے، پر $!C$ تے~$!A$ وچ پابند اے۔
\end{ex}

\begin{defn}[جملہ]
!!^a{formula}~$!A$، \article{sentence}
\emph{!!{sentence}} اُدوں تے صرف اُدوں اے جدوں
اوہدے وچ !!{variable}s دیاں کوئی آزاد آمدواں نہ ہون۔
\end{defn}

% منجمد ماخذ دا غیر مطبوعہ اداراتی نوٹ: ایتھے ہور مثالاں شامل کرن دا اشارہ اے، پر اصل وچ اوہ مثالاں موجود نہیں۔

\end{document}
